\section{Deployment in Azerbaijan}
\label{sec:deployment}

Azerbaijan has no coverage in Sen1Floods11, which motivated a local
mapping prototype. This section describes the prototype, which uses a
different segmentation model and has different evidence from the
event-held-out study.

\subsection{Production Model}

The prototype uses SegFormer-B4 \cite{xie2021segformer} with the
\texttt{nvidia/mit-b4} backbone and ImageNet initialization. It follows
the two-stage training scheme of Sec.~\ref{sec:training}. The logged best
weak-pretraining validation IoU was 0.658. On the hand split the selected
checkpoint reached a validation IoU of 0.7024, test IoU of 0.7027 and
test F1 of 0.8254. Checkpoints within each run were selected by
validation IoU. The historical choice of the default deployment
variant also used test IoU. The
\href{https://github.com/AliAgabalayev/Flood-Damage-Detection/commit/5e727aaed7b5995c6edeec33e90c4e8a87718c8f}{configuration history}
records that two B4 variants had higher validation scores, but the
default B4 beat both variants and the DeepLabV3+ generalist on test
IoU. Table~\ref{tab:model-comparison} records those benchmark runs for
the default and the same generalist as in Table~\ref{tab:two-stage}.
The test contrast is descriptive after model selection. It provides
no accuracy estimate for the Azerbaijan scenes.

\begin{table}[!htbp]
\caption{Logged results for two seed-42 runs on the standard
Sen1Floods11 hand split. These are raw water segmentation scores.}
\label{tab:model-comparison}
\centering\small
\begin{tabular}{@{}lrr@{}}
\toprule
Model & Val. IoU & Test IoU \\
\midrule
DeepLabV3+/ResNet-50 & 0.6712 & 0.6766 \\
SegFormer-B4 & 0.7024 & 0.7027 \\
\bottomrule
\end{tabular}

\end{table}

\subsection{Viewer and API}
\label{sec:api}

The web viewer covers five fixed sites with precomputed assets. It
offers layers for the map mask, recurrent water, model probabilities
and raw SAR, and exports PNG or GeoTIFF where available. A separate
upload page sends a GeoTIFF to the FastAPI \texttt{POST /predict}
endpoint. The model loads once at startup. Inference uses $512\times512$
tiles, 64-pixel overlap, feathered probability stitching and a threshold
of $\geq0.5$.

The revised upload path accepts north-up WGS~84 GeoTIFFs with at least
two bands. It aligns configured JRC and DEM data to the upload grid
and uses the nominal look geometry in Sec.~\ref{sec:fusion}. It
reports partial coverage and rejects requests with no jointly usable
reference coverage. On accepted uploads, missing DEM pixels and
their one-pixel boundary are excluded, while missing JRC values
cause no JRC subtraction. Area, scene percentage and the PNG overlay
come from the same final mask. The percentage uses the full raster as
its denominator.

On 25 November 2026, we sent the original Baku and Sabirabad
GeoTIFFs to a local server with the corrected terrain rule. Each returned HTTP 200 with partial
reference coverage. We decoded the overlay's alpha mask and compared
it pixel by pixel with a fresh offline CLI GeoTIFF using the same
checkpoint, inference settings and reference tiles
(Table~\ref{tab:api-replay}). Both overlays exactly matched the offline
masks. We checked the response bounds and recomputed the map statistics
from the decoded overlays. The final mapped areas were approximately
3.15~km$^2$ in Baku and 2.04~km$^2$ in Sabirabad.

\begin{table}[!htbp]
\caption{Corrected HTTP and CLI replay on the same raster grid.
HTTP and CLI water columns count positive pixels. The last two columns
count flips relative to the previous implementation, for terrain
exclusions and final water separately.}
\label{tab:api-replay}
\centering\small
\begin{tabular}{@{}lrrrrr@{}}
\toprule
Site & HTTP water & CLI water & API--CLI diff. & Terrain flips & Water flips \\
\midrule
Baku & 41,321 & 41,321 & 0 & 892 & 2 \\
Sabirabad & 26,668 & 26,668 & 0 & 0 & 0 \\
\bottomrule
\end{tabular}

\end{table}

\subsection{Scene Provenance}

The project records Baku on 7 April 2026 and Sabirabad on 30 June 2026
as Sentinel-1 GRD IW acquisitions with VV and VH, retrieved using
Google Earth Engine \cite{gorelick2017}. The archived rasters lack
original Sentinel product IDs and overpass times. The Baku City
Executive Power reported heavy rainfall and suburban flooding on
6 April \cite{baku2026}. This provides context for the acquisition.
The records examined for Sabirabad did not confirm a matching event.

\subsection{Reference Labels}

We found no independently labeled inundation extent aligned in time
and grid to either acquisition. Earlier before/after Baku masks were
on a different grid and lacked reliable acquisition and labeling
provenance. The municipal report concerns impacts at city scale and
does not delineate inundation on the radar grid. Both scenes are
treated as unlabeled mapping examples.

\subsection{Effect of Post-Processing}
\label{sec:map-effects}

Table~\ref{tab:deployment} shows the historical offline products and
the CPU replay with the archived exclusion masks. Each replay differs
from the corresponding archived final by one positive pixel. On these
scene grids the approximately mapped area changes from 73.08 to
3.15~km$^2$ in Baku and 8.94 to 2.04~km$^2$ in Sabirabad. The removed
fraction of raw positives is 95.7\% and 77.2\%, respectively. It measures
the effect of the filters on the prediction.

\begin{table}[!htbp]
\caption{Historical offline map products and replay using the archived
terrain rule. All values are positive pixel counts. The corrected replay
is reported separately in Table~\ref{tab:api-replay}.}
\label{tab:deployment}
\centering\small
\begin{tabular}{@{}llrrrr@{}}
\toprule
Site & Date & Raw & After JRC & Replay final & Archived final \\
\midrule
Baku & 2026-04-07 & 960,106 & 165,158 & 41,319 & 41,320 \\
Sabirabad & 2026-06-30 & 116,713 & 44,785 & 26,668 & 26,669 \\
\bottomrule
\end{tabular}

\end{table}

The archived terrain masks use the earlier sign convention, which
reversed the terrain projection in Eq.~(\ref{eq:local-incidence}).
SRTM1 at the nominal geometry reproduces those historical masks.
Dilated DEM gaps account for 1,240,996 of 1,241,953 archived flags in
Baku and all 1,826,806 in Sabirabad. A zero in the archived final
mask may indicate an exclusion or a negative prediction.

The corrected rule leaves these gap counts and the JRC masks
unchanged. In Baku it adds 47 and removes 845 terrain exclusions
outside the gaps, changing 892 terrain pixels. Only two final pixels
change relative to the prior inference on the same inputs, both
returning to the water mask. In Sabirabad, both masks are unchanged.
Relative to the archived final products, the corrected water maps
differ at three pixels in Baku and one in Sabirabad.


\subsection{Limitations}

Each Sen1Floods11 group represents a particular flood event and scene
context. The study measures transfer on the surviving models and
observed chips. The prototype uses a single nominal incidence angle
and local reference coverage, which strongly affects the resulting
map. The two unlabeled scenes leave false alarms, missed water and
the accuracy effect of filtering unmeasured. The reference tiles and
checkpoint used in local replay are listed in the reproducibility
package.
